\documentclass[../../../../SoftwareRequirementsSpecification.tex]{subfiles}
\begin{document}
% Richiede le macro definite in src/macros/useCaseMacros.tex
% (incluse nel preambolo di SoftwareRequirementsSpecification.tex)

\ucstart
\begin{tabularx}{\textwidth}{|l|X|}
  \hline
  \ucid{PT-09} \\ \hline
  \ucname{Personalizzazione di una Scheda Assegnata} \\ \hline
  \ucactor{PT} \\ \hline
  \ucdescription{Il PT interviene sugli allenamenti di
    una scheda che un atleta sta già seguendo, modificandoli solo
    per quell'atleta.} \\ \hline
  \ucprecond{Il PT deve essere autenticato e
    l'atleta scelto deve avere almeno una scheda assegnata in
    corso.} \\ \hline
  \ucpostcond{Gli allenamenti prescritti a
    quell'atleta per la settimana scelta risultano modificati, senza
    che cambino né la scheda né gli allenamenti degli altri
    atleti che la seguono.} \\ \hline
  \ucmainflow{
    \item Il PT si trova sulla pagina di dettaglio di un atleta,
      raggiunta tramite il caso d'uso PT-08.
    \item Il PT preme su una scheda dell'elenco delle schede
      assegnate.
    \item Il sistema mostra il dettaglio dell'assegnazione, che
      riporta il nome della scheda, il periodo di validità e
      l'elenco delle settimane di allenamento, con evidenziata la
      settimana corrente e segnalate le settimane che contengono
      personalizzazioni.
    \item Il PT preme su una settimana, scelta fra quella corrente e
      le successive (vedi Nota 2).
    \item Il sistema mostra le sessioni di allenamento della
      settimana con i valori previsti per quell'atleta (vedi
      Nota 1).
    \item Il PT preme su un esercizio.
    \item Il sistema mostra il form con i valori correnti
      dell'esercizio:
      \begin{itemize}
        \item peso;
        \item numero di serie/ripetizioni;
        \item tempo di recupero tra una serie e l'altra (in
          secondi).
      \end{itemize}
    \item Il PT modifica i valori che vuole ritoccare e indica se
      estendere la modifica anche alle settimane successive del
      blocco mensile (vedi Nota 5).
    \item Il PT preme il tasto "Salva".
    \item Il sistema registra la personalizzazione per la settimana
      scelta e, se il PT ne ha chiesto l'estensione, la registra
      anche per ciascuna settimana rimanente del blocco, ricalcolando
      i valori con la progressione dell'esercizio a partire dal nuovo
      valore inserito.
    \item Il sistema invia una mail di notifica all'Atleta per
      informarlo che la sua scheda è cambiata.
  } \\ \hline
  \ucaltflow{
    \textbf{2a. Nessuna scheda assegnata} \newline
    - L'atleta non ha schede assegnate: il sistema mostra la
    sezione vuota. \newline
    - Il flusso termina. \newline
    \textbf{4a. Scheda terminata} \newline
    - Tutte le settimane della scheda sono trascorse e nessuna è
    quindi modificabile. \newline
    - Il sistema mostra un messaggio ("Scheda terminata: non ci sono
    settimane modificabili"). \newline
    - Il flusso termina. \newline
    \textbf{7a. Ripristino del valore originale} \newline
    - Se l'esercizio è stato personalizzato, il sistema mostra
    accanto ai valori il tasto "Ripristina". \newline
    - Il PT lo preme e il sistema riporta l'esercizio ai valori
    previsti dalla scheda. \newline
    - Se lo stesso esercizio è personalizzato anche nelle settimane
    successive del blocco, il sistema chiede se ripristinare anche
    quelle. \newline
    - Il flusso riprende dal punto 5. \newline
    \textbf{8a. Salto di un esercizio} \newline
    - Invece di ritoccare i valori il PT preme il tasto "Salta
    esercizio", così che l'atleta non lo svolga in quella
    settimana. \newline
    - Il flusso riprende dal punto 10. \newline
    \textbf{8b. Sostituzione di un esercizio} \newline
    - Invece di ritoccare i valori il PT preme il tasto
    "Sostituisci". \newline
    - Il sistema mostra la lista degli esercizi, come al passo 6 del
    caso d'uso PT-04. \newline
    - Il PT sceglie l'esercizio sostitutivo e ne configura i valori.
    \newline
    - Il flusso riprende dal punto 10.
  } \\ \hline
  \ucnote{
    \textbf{Nota 1:} i valori mostrati al passo 5 non sono
    memorizzati: il sistema li calcola a partire dagli allenamenti
    configurati nel blocco mensile (PT-04), dalla progressione di
    ciascun esercizio, dai pesi inseriti per l'atleta al momento
    dell'assegnazione (PT-05) e dalle personalizzazioni già
    presenti. La modifica è registrata come scostamento sulla
    singola assegnazione: la scheda resta invariata e gli altri
    atleti che la seguono non sono toccati.\newline
    \textbf{Nota 2:} sono modificabili solo la settimana corrente e
    le successive. Le settimane già trascorse sono in sola lettura,
    perché l'atleta si è già allenato su quanto vi era
    prescritto.\newline
    \textbf{Nota 3:} lo statement prescrive che la tipologia di
    esercizi in una scheda vari solo da mese a mese. Il vincolo è
    rispettato sulla scheda, dove la tipologia è fissata per
    l'intero blocco mensile dagli allenamenti configurati in PT-04.
    Il salto e la sostituzione di un esercizio (flussi 8a e 8b) non
    modificano la scheda: sono eccezioni riferite a un singolo atleta
    e a una singola settimana, che rispondono a situazioni
    contingenti quali un infortunio o l'indisponibilità di un
    macchinario.\newline
    \textbf{Nota 4:} la personalizzazione non consente di aggiungere
    esercizi non previsti, di cambiare il numero o la composizione
    delle sessioni della settimana, di cambiare la durata della
    scheda, né di cambiare la progressione assegnata a un
    esercizio, che è parte della scheda e vale per tutti gli atleti
    che la seguono. Per interventi di questa natura il PT deve
    duplicare la scheda e assegnare la copia.\newline
    \textbf{Nota 5:} l'estensione del passo 8 si ferma alla fine del
    blocco mensile corrente, perché i blocchi sono indipendenti e
    la progressione riparte a ogni blocco. Se l'esercizio non ha una
    progressione, il valore inserito è ripetuto identico sulle
    settimane rimanenti.\newline
    \textbf{Nota 6:} la personalizzazione riguarda ciò che è
    prescritto all'atleta e non altera in alcun modo gli allenamenti
    che l'atleta ha già svolto e registrato.
  } \\ \hline
\end{tabularx}
\end{document}
